
\subsection{Invariant Generation for Loops}
Algorithm~\supref{alg:loopinv-gen} (\appref{app:loop-analysis}) takes a loop entry $p_{\textit{loop}}$, its symbolic precondition $P$, and an LLM to synthesize an invariant; other inputs are as defined above.

The procedure has three steps. First, the LLM reads the outer loop body and answers guiding questions; this supports later synthesis but is never emitted as an annotation (Line~1). Second, $\mathrm{LoopAnalysis}$ processes inner loops and symbolically executes the body, returning unchanged variables $\textit{uV}$ and non-inductive variables $\textit{nV}$ while restoring the loop-stack depth (Lines~2--4). Third, $\mathrm{LoopInvTemGen}$ builds templates from that analysis, $\mathrm{OuterLoopLLMCall}$ fills them, and $\mathrm{Refine}$ repairs verification failures (Lines~5--9; Section~\ref{sec:refinement}).

\rev{$\mathrm{LoopAnalysis}$ does this by symbolically executing a
single-iteration unfolding of the loop body and classifying $\textit{uV}$ and
$\textit{nV}$ from the resulting postcondition; nested loops, for which no
well-defined precondition exists at the enclosing loop boundary, are discharged
bottom-up by direct LLM synthesis ($\mathrm{NestedLoopLLMCall}$) without
templates. The full procedure is given in \appref{app:loop-analysis} (Algorithm~\supref{alg:loop-analysis}).}
